% SİFPGA Doğrulama VIP Seti - Tasarım Dokümanı (taslak)
% Overleaf: pdfLaTeX ile derlenir (Menü -> Compiler -> pdfLaTeX).
\documentclass[11pt,a4paper]{article}

\usepackage[T1]{fontenc}
\usepackage[utf8]{inputenc}
\usepackage[turkish]{babel}
\usepackage{lmodern}
\usepackage[margin=2.5cm]{geometry}
\usepackage{booktabs}
\usepackage{longtable}
\usepackage{array}
\usepackage{enumitem}
\usepackage{xcolor}
\usepackage{tikz}
\usetikzlibrary{arrows.meta,positioning,calc,fit}
\usepackage[hidelinks]{hyperref}
\usepackage{fancyhdr}
\usepackage{textcomp}
\usepackage{amssymb}

\setlength{\parindent}{0pt}
\setlength{\parskip}{0.6em}
\renewcommand{\arraystretch}{1.2}

\pagestyle{fancy}
\fancyhf{}
\lhead{SİFPGA Doğrulama VIP Seti}
\rhead{Tasarım Dokümanı (taslak)}
\cfoot{\thepage}

% Kısayollar
\newcommand{\us}{\,\textmu{}s}
\newcommand{\kod}[1]{\texttt{#1}}
\newcommand{\acik}[1]{\textcolor{red!60!black}{\textbf{[AÇIK]} #1}}

% TikZ stilleri
\tikzset{
  kutu/.style={draw=black!60, rounded corners=3pt, align=center,
               minimum height=1.1cm, inner sep=4pt, font=\small},
  ana/.style={kutu, draw=blue!70!black, fill=blue!8, very thick},
  ok/.style={-{Stealth[length=2.2mm]}, thick, black!70},
  etiket/.style={font=\scriptsize, text=black!70}
}

\title{\textbf{SİFPGA Doğrulama VIP Seti}\\[0.3em]
       \large Tasarım Dokümanı (taslak)}
\author{}
\date{\today}

\begin{document}
\maketitle
\tableofcontents
\newpage

%---------------------------------------------------------------
\section{Amaç ve kapsam}

Bu VIP seti, radar sinyal işleme FPGA'sını (SİFPGA) gerçek RTL ile birlikte
Vivado/xsim ortamında doğrular. CPU'dan hüzme parametreleri gelir, iki SHOB
kartından Aurora ile veri gelir, SİFPGA sonucu PCIe üzerinden RSİ belleğine
yazar ve MSI kesmesi üretir.

\paragraph{Kapsam içi}
\begin{itemize}[nosep]
  \item Hüzme turu akışı: HB planlama, hüzme parametre mesajları (HUZPAR),
        SHOB hüzme verisi, RSİ çıktısı.
  \item Çoklu hüzme modu ve tur başına rastgele parametre (M, darbe, örnek sayısı).
  \item RSİ çıktısının deterministik veri deseniyle sayısal doğrulaması.
  \item SİFPGA'dan gelen MSI1--MSI4 kesmelerinin kaydı.
\end{itemize}

\paragraph{Kapsam dışı}
\begin{itemize}[nosep]
  \item Tasarımcıya ait \kod{Top.vhd} ve \kod{PcieTop.vhd} dosyaları değiştirilmez.
  \item Hata enjeksiyonu (MSI3/MSI4 üretimi) ve \kod{0x00000003} (Bypass) RSİ
        mesajı henüz modellenmedi.
  \item \kod{1DEA001B} (hüzme veri durumu) ve \kod{1DEADEAD} (hata) mesajlarının
        çözümleyicisi henüz yok; anlamları tasarımcıdan teyit bekliyor.
\end{itemize}

%---------------------------------------------------------------
\section{Mimari genel bakış}

Senaryo okuyucu tek bir metin dosyasını okur ve CPU/PCIe, iki SHOB ve HB
sürücülerini yönlendirir. SİFPGA'nın çıktısı PCIe RQ üzerinden RSİ modeline
gider; model MSI'ları da kaydeder.

\begin{figure}[h]
\centering
\begin{tikzpicture}[node distance=0.5cm and 1.3cm]
  \node[kutu, text width=2.6cm] (okuyucu) {\textbf{Senaryo okuyucu}\\ metin dosyası,\\ \kod{\$fgetc} ile akış};
  \node[kutu, text width=3.0cm, right=1.3cm of okuyucu, yshift=2.4cm] (cpu)
        {\textbf{CPU / PCIe VIP}\\ \kod{1DEA0001}, \kod{1DEA0003},\\ IUB, SHOB hüzpar};
  \node[kutu, text width=3.0cm, below=0.5cm of cpu] (s1)
        {\textbf{SHOB-1 VIP}\\ Aurora hüzme verisi\\ HB'den 5\us{} sonra};
  \node[kutu, text width=3.0cm, below=0.5cm of s1] (s2)
        {\textbf{SHOB-2 VIP}\\ Aurora hüzme verisi\\ HB'den 5\us{} sonra};
  \node[kutu, text width=3.0cm, below=0.5cm of s2] (hb)
        {\textbf{HB üreteci}\\ planlı HB darbesi\\ \kod{77777777} komutu};
  \node[ana, minimum height=7.6cm, text width=1.6cm, right=1.6cm of s1, yshift=-1.0cm] (sif)
        {\textbf{SİFPGA}\\ doğrulanan RTL};
  \node[kutu, text width=2.4cm, right=1.6cm of sif] (rsi)
        {\textbf{RSİ modeli}\\ ID ile eşleştirir\\ deseni denetler\\ MSI'ları kaydeder};

  \foreach \n in {cpu,s1,s2,hb}{
    \draw[ok] (okuyucu.east) -- ++(0.5,0) |- (\n.west);
  }
  \draw[ok] (cpu.east) -- node[etiket,above]{PCIe CQ} (cpu.east -| sif.west);
  \draw[ok] (s1.east)  -- node[etiket,above]{Aurora}  (s1.east -| sif.west);
  \draw[ok] (s2.east)  -- node[etiket,above]{Aurora}  (s2.east -| sif.west);
  \draw[ok] (hb.east)  -- node[etiket,above]{HB}      (hb.east -| sif.west);
  \draw[ok] (sif.east) -- node[etiket,above]{PCIe RQ} node[etiket,below]{MSI} (rsi.west);
\end{tikzpicture}
\caption{Doğrulama ortamı: dört sürücü, SİFPGA ve RSİ modeli.}
\end{figure}

%---------------------------------------------------------------
\section{Arayüzler}

Tüm CPU ve SHOB trafiği SİFPGA'nın PCIe uç noktasına BAR numarasıyla yönlenir;
SHOB hüzme verisi ise Aurora ile ayrı gelir.

\begin{table}[h]
\centering
\begin{tabular}{@{}lll@{}}
\toprule
Arayüz & Yön & Açıklama \\
\midrule
PCIe CQ (AXI-S, 256 bit) & VIP $\to$ SİFPGA & CPU, IUB ve SHOB hüzpar yazımları \\
PCIe RQ (AXI-S, 256 bit) & SİFPGA $\to$ VIP & RSİ bellek yazımları, MSI \\
Aurora (SHOB-1, SHOB-2) & SHOB VIP $\to$ SİFPGA & Hüzme verisi, HB'den sonra gecikmeli \\
HB & VIP $\to$ SİFPGA & Hüzme başı darbesi, \kod{77777777} ile planlanır \\
\bottomrule
\end{tabular}
\caption{Arayüzler.}
\end{table}

\begin{table}[h]
\centering
\begin{tabular}{@{}ll@{}}
\toprule
BAR & Hedef \\
\midrule
BAR0 & SİFPGA (CPU mesajları: \kod{1DEA0001}, \kod{1DEA0003}) \\
BAR1 & SHOB-1 hüzpar (\kod{0000E001}) \\
BAR2 & SHOB-2 hüzpar (\kod{0000E001}) \\
BAR3 & IUB mesajı (\kod{00000001}) \\
\bottomrule
\end{tabular}
\caption{BAR haritası.}
\end{table}

\begin{table}[h]
\centering
\begin{tabular}{@{}lll@{}}
\toprule
MSI & Anlamı & Kaynak \\
\midrule
MSI1 & RSİ mesajı yazıldı & SİFPGA \\
MSI2 & RSİ yazımı başarısız & SİFPGA \\
MSI3 & SHOB$\to$SİFPGA kanal verisi eksik & SİFPGA \\
MSI4 & SHOB verisi alınamadı & SİFPGA \\
\bottomrule
\end{tabular}
\caption{MSI kesmeleri. MSI'ları SİFPGA üretir; RSİ modeli yalnızca alıp loga yazar.}
\end{table}

%---------------------------------------------------------------
\section{Senaryo dili ve üretici}

Tüm girdi tek bir metin dosyasıdır; her satır bir PCIe yazımı ya da bir
zamanlama komutudur.

\paragraph{Satır biçimi.}
\kod{ADRES [BAR<n>] MSGID LEN DATA...} Veri baytları küçük-endian dizilir.
32 DWORD'den uzun mesajlar otomatik olarak 32 DWORD'lük TLP'lere bölünür.

\begin{table}[h]
\centering
\begin{tabular}{@{}ll@{}}
\toprule
Komut & Anlamı \\
\midrule
\kod{77777777 <ns>} & Bir sonraki HB'yi \kod{<ns>} sonraya planla \\
\kod{99999999} & Planlı HB'yi bekle \\
\kod{88888888 <ns>} & \kod{<ns>} bekle \\
\bottomrule
\end{tabular}
\caption{Zamanlama komutları.}
\end{table}

\paragraph{Okuyucu.}
Senaryo dosyası \kod{\$fgetc} ile akış olarak okunur. 550~KB'lık satırları
\kod{\$fgets} ile belleğe almak xsim'de bellek ve süre sorunu çıkardığından
v1.0.8'de bu yöntem bırakıldı.

\paragraph{Üreticiler (\kod{tools/}).}
\begin{itemize}[nosep]
  \item \kod{gen\_huzme\_rounds.py}: sabit parametreli turlar için yardımcı
        fonksiyonlar (\kod{le}, \kod{line}, \kod{shob\_huzpar}, \kod{iub\_huzpar},
        \kod{sifpga\_huzpar}).
  \item \kod{gen\_huzme\_random.py}: tur başına rastgele M, darbe, örnek ve RSİ
        başına hüzme dağılımı üretir. Seçenekler: \kod{--count}, \kod{--seed},
        \kod{--m-max} (12), \kod{--n-max} (64), \kod{--darbe-max}, \kod{--ornek-max},
        \kod{--gaps-us} (20,30), \kod{--huzme-id}.
  \item Tur sırası: \kod{77777777}, \kod{1DEA0003}, IUB (BAR3), SHOB-1 (BAR1),
        SHOB-2 (BAR2), \kod{99999999}, \kod{88888888 <boşluk>}. Hüzme ID'si her
        turda 1 artar.
\end{itemize}

%---------------------------------------------------------------
\section{Hüzme turu akışı}

\begin{figure}[h]
\centering
\begin{tikzpicture}[node distance=0.55cm]
  \node[kutu, text width=10.5cm] (a) {\textbf{Planlama} \quad \kod{77777777} komutu HB'yi planlar};
  \node[kutu, text width=10.5cm, below=of a] (b) {\textbf{HUZPAR} \quad \kod{1DEA0003}, IUB, SHOB-1, SHOB-2 PCIe ile gider};
  \node[ana, text width=10.5cm, below=of b] (c) {\textbf{HB} \quad HB darbesi (periyot en az 60\us)};
  \node[kutu, text width=10.5cm, below=of c] (d) {\textbf{SHOB verisi} \quad Her iki SHOB HB'den 5\us{} sonra Aurora ile basar};
  \node[kutu, text width=10.5cm, below=of d] (e) {\textbf{RSİ yazımı} \quad SİFPGA belleğe yazar, MSI1 gelir (MSI3 bazı turlarda)};
  \node[kutu, text width=10.5cm, below=of e] (f) {\textbf{Bekleme} \quad \kod{88888888}: 20/30/100\us, sonra yeni tur};
  \foreach \x/\y in {a/b,b/c,c/d,d/e,e/f}{ \draw[ok] (\x.south) -- (\y.north); }
  \draw[ok] (f.east) -- ++(0.7,0) |- node[etiket,pos=0.25,right]{yeni tur} (a.east);
\end{tikzpicture}
\caption{Bir hüzme turunun adımları (zamana ölçekli değil).}
\end{figure}

MSI3, bazı turlarda MSI1'den önce gelir ve veriyi bozmaz. Boşluk 20/30\us{}
olduğunda büyük N'li turlardan sonra bazı turlar başarısız olur; 100\us{}'de
olmaz (bkz.\ Bölüm~\ref{sec:gozlem}).

%---------------------------------------------------------------
\section{Mesaj biçimleri}

Tüm alanlar küçük-endian'dır.

\subsection{CPU $\to$ SİFPGA}

\begin{longtable}{@{}p{2.6cm}p{1.6cm}p{10.2cm}@{}}
\toprule
Mesaj & Boyut & İçerik \\
\midrule
\endhead
\kod{1DEA0001} & 124 B & 4 tampon tanımı (taban, boyut, işaretçi; 8~B'lık alanlar),
96..119 arası 3 RSİ başlangıç adresi, RSİ sayısı, kesme aralığı (\textmu{}s) \\
\kod{1DEA0003} & 56 B & Hüzme parametresi (aşağıda) \\
IUB \kod{00000001} & BAR3 & IUB hüzpar'ı \\
SHOB \kod{0000E001} & BAR1/2 & SHOB hüzpar'ı \\
\bottomrule
\end{longtable}

\paragraph{\kod{1DEA0003} alanları (56 B).}

\begin{table}[h]
\centering
\begin{tabular}{@{}ll@{}}
\toprule
Bayt & Alan \\
\midrule
0--23 & 3 adet 8~B RSİ ofseti \\
24--27 & Hüzme ID \\
28--29 & M (toplam hüzme sayısı) \\
30--31 & Darbe sayısı \\
32--33 & Örnek sayısı \\
34 & SHOB kullanım (03) \\
35 & Mod (0 = HO) \\
36--38 & $n_1, n_2, n_3$: RSİ başına hüzme sayısı (toplam $=M$) \\
39--55 & Ayrılmış \\
\bottomrule
\end{tabular}
\caption{\kod{1DEA0003} yerleşimi.}
\end{table}

$N = \text{darbe} \times \text{örnek}$. Bir RSİ mesajındaki kanal sayısı $K$,
ilgili $n$ değerine eşittir.

\subsection{SHOB Aurora verisi}

Kimlik 2, ardından $\mathrm{LEN} = 56 + M\cdot N\cdot 8 + 960 + 16 + 16$ bayt:
başlık, örnekler (örnek-ana sıra, int32 I ve Q), kuyruk alanları.

\subsection{RSİ mesajı (SİFPGA $\to$ bellek)}

\begin{table}[h]
\centering
\begin{tabular}{@{}ll@{}}
\toprule
Alan & Boyut \\
\midrule
ID ($=2$), LEN & 4~B + 4~B \\
Başlık & 88~B \\
Kanal başına $N$ adet I, sonra $N$ adet Q (float32) & $8N$~B her kanal için \\
GAC SCIT & 1920~B \\
SHOB SCIT & 32~B \\
SİBFPGA SCIT & 16~B \\
Alt bilgi (\kod{AAAAAAAA BBBBBBBB CCCCCCCC DDDDDDDD}) & 16~B \\
\bottomrule
\end{tabular}
\caption{RSİ mesaj yerleşimi.}
\end{table}

Sabit kısım 2080~B'tır:
\[
  K = \frac{\text{toplam bayt} - 2080}{8\,N}.
\]

\subsection{SİFPGA $\to$ CPU (gözlenen)}

Anlamlar kısmen tahmindir.

\begin{longtable}{@{}p{2.6cm}p{12.3cm}@{}}
\toprule
Mesaj & Gözlem \\
\midrule
\endhead
\kod{1DEA0002} & \kod{1DEA0001} onayı \\
\kod{1DEA0004} & \kod{1DEA0003} olumlu cevabı: hüzme ID + durum 0 \\
\kod{1DEA001B} & 16 B: ID, A, 0, B. Normal turda tümü sıfır; MSI3'lü turda iki rapor
(ID$-$1 ile A$=1$, sonra ID ile A$=0$), B$=3$ \\
\kod{1DEADEAD} & 8 B: ID + hata kodu. 0x0F, ID mod 4 $=3$ olan her turda geliyor;
anlamı bilinmiyor \\
\bottomrule
\end{longtable}

%---------------------------------------------------------------
\section{Veri deseni ve beklenen sonuç}

SHOB verisi örnek ve hüzme indeksinden üretilen deterministik bir desendir;
SİFPGA iki SHOB'u topladığından RSİ'de beklenen değer üç katıdır.

\begin{itemize}[nosep]
  \item Kod: $\mathrm{kod} = ((m+1) \ll 15)\,|\,(n \mathbin{\&} \mathtt{0x7FFF})$;
        $m$ hüzme, $n$ örnek indeksi.
  \item SHOB-1: $I = +1\cdot\mathrm{kod}$, $Q = -1\cdot\mathrm{kod}$.
        SHOB-2: $I = +2\cdot\mathrm{kod}$, $Q = -2\cdot\mathrm{kod}$.
  \item Toplam: $I = 3\,\mathrm{kod}$, $Q = -3\,\mathrm{kod}$.
        Örnek: $m=0$ için $I = 98304 + 3n$.
  \item RSİ değerleri float32'dir; değerler $2^{24}$'ten küçük olduğundan
        tam gösterilir.
\end{itemize}

\paragraph{Zamanlama.}
SHOB VIP verisi, HUZPAR'dan sonraki ilk HB'den \kod{DATA\_DELAY\_CLKS} $=625$ clk
(5\us) sonra başlatır. HB, HUZPAR mesajlarından sonra gelmelidir; üretici HB
periyodunu $\max(60\us,\ 20\us + 3\us\times M)$ seçer.

%---------------------------------------------------------------
\section{RSİ modeli ve doğrulama kuralları}

RSİ modeli (\kod{sim\_model\_rsi.sv}, \kod{rsi\_pkg.sv}) her gelen mesajı hüzme
ID'siyle bir tura eşler, $K$'yı bayt sayısından hesaplar ve her örneği beklenen
desenle karşılaştırır.

\begin{enumerate}[nosep]
  \item \textbf{Eşleştirme:} Başlıktaki hüzme ID'sine göre tur seçilir. Beklenen
        sıradan ileri atlanırsa aradaki turlar ``VERI GELMEDI'' sayılır ve loga
        \kod{ESLESME} satırı yazılır. Sıralı eşleştirme, reddedilen bir turdan
        sonra yanlış zincirleme hataya yol açıyordu.
  \item \textbf{K doğrulaması:} $K=(\text{bayt}-2080)/(8N)$. \kod{1DEA0003}'teki
        $n$ değeriyle uyuşmazsa hata.
  \item \textbf{Boyut doğrulaması:} Gelen bayt sayısı $8+\mathrm{LEN}$ ile
        karşılaştırılır; fark \kod{LENFARK} olarak raporlanır. Gelen boyut önceki
        turla aynıysa teşhis satırı bunu belirtir.
  \item \textbf{Sayısal kontrol:} Örnek başına I ve Q, toplam kuralıyla
        ($\times 3$) karşılaştırılır; GAC/SHOB/SİBFPGA SCIT ve alt bilgi ayrıca
        kontrol edilir.
  \item \textbf{MSI kaydı:} MSI3 sayısı bir sonraki turun \kod{TESHIS} satırında
        gösterilir.
\end{enumerate}

Tur sonucu üç değerden biridir: DOĞRU, YANLIŞ, VERİ GELMEDİ. Sonunda
\kod{rsi\_tur\_ozeti} dosyası tüm turları özetler.

%---------------------------------------------------------------
\section{Gözlenen davranışlar ve sınırlar}
\label{sec:gozlem}

Aşağıdakiler gerçek SİFPGA ile yapılan koşulardan çıkarılan gözlemlerdir;
tasarımcı teyidi bekliyor.

\begin{longtable}{@{}p{3.2cm}p{11.7cm}@{}}
\toprule
Konu & Gözlem \\
\midrule
\endhead
$M=16$ & Hep reddedildi (NAK, hata kodu 03: hatalı hüzme sayısı parametresi) \\
$N$ sınırı & $M\le 8$'de $N\ge 80$ reddedildi; $M=12$'de $N=112$'ye kadar kabul edildi \\
Reddedilen tur & RSİ'ye yine yazım yapılıyor, boyut bir önceki turunki \\
Turlar arası süre & 20/30\us{} boşluk, $N\ge 56$ turundan sonra 4/100 tur başarısız
(2 kayıp, 2 eski boyutlu). 100\us{} boşlukta 100/100 doğru \\
MSI3 & 100\us{} koşusunda 99 turun 61'inde var, boşluktan bağımsız. HB döngüsünde
sabit fazda geliyor, MSI1'den 6--9\us{} önce. MSI3'lü turlar yine de doğru veriyle
bitiyor \\
\kod{1DEA001B} & MSI3 olan turlarla birebir eşleşiyor \\
\kod{1DEADEAD} 0x0F & ID mod 4 $=3$ olan 25 turun hepsinde, MSI3'ten ve boyuttan
bağımsız; RSİ verisi doğru \\
\bottomrule
\end{longtable}

$n_1/n_2/n_3$ dağılımının kabul için bir kuralı olmadığı görüldü. MSI3 hakkında
``SHOB-2 verisi kontrol anına yetişmiyor'' tahmini var, test edilmedi.

%---------------------------------------------------------------
\section{Açık sorular ve varsayımlar}

\begin{itemize}[label=$\square$]
  \item \kod{1DEADEAD} hata kodu 0x0F ne demek ve neden ID mod 4 $=3$ olan
        turlarda geliyor? Hangi mesaja cevap veriyor?
  \item \kod{1DEA001B} alanları A ve B ne? Neden MSI3'lü turda önceki ID için de
        rapor geliyor?
  \item MSI3 hangi koşulda ve HB'den kaç \textmu{}s sonra üretiliyor? SHOB verisinin
        HB'ye göre en geç başlama zamanı nedir?
  \item Kabul sınırları: M üst sınırı, N üst sınırı ve hata kodu 03'ün kesin koşulu.
  \item Reddedilen turda neden önceki turun boyutuyla yazım yapılıyor?
  \item Büyük bir turdan sonra yeni tur kabul etmek için asgari süre nedir?
        (20--30\us{} yetmiyor, 100\us{} yetiyor.)
  \item IUB cevabındaki ``hüzme başı zamanı'' birimi ve çözünürlüğü (ns, 5~ns?).
  \item RSİ tarafı: eksik ya da eski ID'li turlar gerçek RSİ'de nasıl işleniyor?
        \kod{0x00000003} (Bypass) mesajı ne zaman kullanılıyor?
\end{itemize}

\paragraph{Varsayımlar.}
M ve N sınırları gözlemden çıkarılmıştır. SHOB üzerindeki toplam kuralı ($\times3$)
tüm turlarda tutuyor, ama tasarım belgesinden teyit edilmedi.

%---------------------------------------------------------------
\section{Çalıştırma ve test}

Set karma Verilog/SystemVerilog olduğundan Vivado xsim 2020.2 ile çalışır; birim
testleri Verilator ile koşulur.

\begin{itemize}[nosep]
  \item \textbf{Sürüm klasörleri:} \kod{verification\_v1.0.7} (dokunulmaz taban),
        \kod{verification\_v1.0.8} (akış okuyucu), \kod{verification\_v1.0.9} (güncel).
  \item \textbf{Verilator:} \kod{test/run\_verilator.sh} 9 testi çalıştırır; her
        v1.0.9 commit'inde geçti.
  \item \textbf{Bellek teşhisi:} \kod{test/xsim\_bellek\_teshis/} xsim'de bellek
        tüketiminin hangi aşamada olduğunu ayırmak için kullanılır.
  \item \textbf{xsim sınırları:} Fonksiyon içinden task çağrılamaz; çıkış
        argümanları ve diziler dikkat ister.
  \item \textbf{Örnek girdiler} (\kod{examples/}):
        \kod{input\_huzme\_100\_rastgele.txt} (20/30\us{} boşluk),
        \kod{input\_huzme\_100\_rastgele\_bosluk100us.txt} (100\us),
        \kod{input\_huzme\_nak\_teshis.txt} (M/N sınır ayrıştırma),
        \kod{input\_huzme\_MAX\_tek\_huzme.txt}.
  \item \textbf{Çıktılar:} RSİ modeli tur ayrıntılarını ve \kod{rsi\_tur\_ozeti}
        özetini yazar. Referans doğrulama için 100\us{} boşluklu dosya önerilir.
\end{itemize}

%---------------------------------------------------------------
\section{Sürüm geçmişi}

\begin{table}[h]
\centering
\begin{tabular}{@{}lp{12cm}@{}}
\toprule
Sürüm & Değişiklik \\
\midrule
v1.0.9 & RSİ modeli hüzme ID'siyle eşleştirir; kayıp tur, LENFARK ve MSI3 teşhisleri;
SİFPGA kabul sınırlarına göre rastgele üretici; 100\us{} boşluklu dosya \\
v1.0.8 & Senaryo dosyası \kod{\$fgetc} ile akış olarak okunur (xsim'de 17--32~GB
bellek sorununu giderir) \\
v1.0.7 & Taban sürüm; değiştirilmeden saklanıyor \\
\bottomrule
\end{tabular}
\caption{Sürüm geçmişi.}
\end{table}

\end{document}
